\documentclass[10pt,a4paper,twocolumn]{article}
\usepackage[utf8]{inputenc}
\usepackage[T1]{fontenc}
\usepackage[italian]{babel}
\usepackage[margin=1.8cm,top=2cm,bottom=2cm]{geometry}
\usepackage{graphicx}
\usepackage{booktabs}
\usepackage{amsmath}
\usepackage{listings}
\usepackage{xcolor}
\usepackage{hyperref}
\usepackage{enumitem}
\usepackage{caption}

\setlength{\columnsep}{0.6cm}
\setlist{noitemsep,topsep=2pt}
\captionsetup{font=small,labelfont=bf}

\lstset{
  basicstyle=\ttfamily\scriptsize,
  keywordstyle=\color{blue!70!black},
  commentstyle=\color{green!50!black},
  breaklines=true,
  frame=single,
  xleftmargin=2mm,
  framexleftmargin=2mm,
  columns=fullflexible,
}

\title{\vspace{-1.5em}\textbf{SPM Modulo 1: Vettorizzazione del kernel di Partition Mapping}\\\small\textit{(Versione italiana --- solo per uso personale, non inclusa nella consegna)}\vspace{-0.5em}}
\author{Gabriele Marsili}
\date{}

\begin{document}
\maketitle
\vspace{-1.5em}

% ===========================================================================
\section{Scelte progettuali}
% ===========================================================================

Dati $N$ chiavi (\texttt{uint64\_t}) e $P$ partizioni (potenza di due), il compito è calcolare $\text{part\_id}[i] = h(\text{keys}[i]) \in [0,P)$ per ogni chiave.

\subsection{Funzione di mapping}

Ho scelto una hash Fibonacci multiply-shift che opera su 32~bit:
$$h(k) = \bigl((k_\text{lo} \oplus k_\text{hi}) \cdot A_{32}\bigr) \gg (32 - \log_2 P)$$
dove $A_{32} = \texttt{0x9E3779B9} = \lfloor 2^{32}/\varphi \rfloor$~\cite{knuth1998}.
Si tratta di un'istanza dello schema di \emph{hashing by multiplication} descritto in~\cite{ferragina2023} (\S\,8.2), dove la costante legata al rapporto aureo $A = (\sqrt{5}-1)/2 \approx 0.618$ è raccomandata per le sue proprietà di dispersione dei bit. L'estrazione tramite right-shift segue la famiglia universale di hash \emph{multiply-add-shift} di~\cite{ferragina2023} (\S\,8.3): $h_{a,b}(k)=(ak+b \bmod 2^w) \operatorname{div} 2^{w-\ell}$, qui semplificata con $b\!=\!0$ e la costante di Fibonacci come~$a$.

Lo XOR delle due metà a 32~bit preserva l'informazione di entrambe le metà prima della moltiplicazione. La costante di Fibonacci distribuisce bene i bit nell'intervallo di output: la distribuzione misurata su $10^8$ chiavi con $P\!=\!256$ dà $\max(\text{count})/\text{atteso} = 1{,}005$.

La ragione principale per lavorare a 32~bit è la compatibilità SIMD: \texttt{\_mm256\_mullo\_epi32} (\texttt{vpmulld}) è una singola istruzione nativa AVX2 che moltiplica 8~valori uint32 in parallelo. Una moltiplicazione a 64~bit in AVX2 richiederebbe 3~istruzioni \texttt{vpmuludq} più shift e addizioni per elemento --- ho verificato che risulta più lenta della baseline scalare in questo regime limitato dalla banda di memoria.

\subsection{Layout dei dati e P}

Chiavi e partition ID di output risiedono in array contigui separati con stride unitario --- il layout più favorevole alla vettorizzazione. La memoria è allineata a 32~byte tramite \texttt{posix\_memalign} per i load allineati AVX2. $P$ è una potenza di due, così il modulo si riduce a un right-shift.

% ===========================================================================
\section{Auto-vettorizzazione}
% ===========================================================================

Il kernel C++ ``plain'' viene compilato da un singolo file sorgente in due binari: uno con \texttt{-fno-tree-vectorize} (baseline), l'altro con \texttt{-O3 -march=native -mavx2 -ftree-vectorize} (autovec). Entrambi usano \texttt{-fopt-info-vec-optimized} e \texttt{-fopt-info-vec-missed} per la diagnostica.

GCC~12.2 riporta la vettorizzazione riuscita del loop del kernel. Le righe rilevanti dal report completo (file \texttt{vec\_report\_optimized.txt}, incluso nei deliverable):

\begin{lstlisting}
src/plain.cpp:33:26: optimized: loop vectorized
  using 32 byte vectors
src/plain.cpp:33:26: optimized: loop vectorized
  using 16 byte vectors
include/common.hpp:77:27: optimized: loop vectorized
  using 32 byte vectors
include/common.hpp:264:19: optimized: loop vectorized
  using 32 byte vectors
\end{lstlisting}

GCC emette due versioni del loop del kernel: una che usa registri AVX2 a 256~bit (\texttt{ymm}, vettori da 32~byte, 8~uint32 per istruzione) e una che usa registri SSE a 128~bit (\texttt{xmm}, vettori da 16~byte, 4~uint32 per istruzione). A runtime viene preso il percorso a 256~bit quando AVX2 è disponibile; la versione a 128~bit serve come fallback su hardware più vecchio. Le righe relative a \texttt{common.hpp} corrispondono ai loop di checksum e generazione chiavi, che GCC vettorizza anch'essi. Il report \texttt{-fopt-info-vec-missed} è risultato vuoto --- GCC non ha trovato alcun loop che abbia tentato di vettorizzare senza riuscirci.

Il loop soddisfa tutte le condizioni: trip count noto, nessuna dipendenza loop-carried, nessuna chiamata a funzione nel corpo, \texttt{\_\_restrict\_\_} per escludere l'aliasing, e accesso con stride unitario.

% ===========================================================================
\section{Intrinsics AVX2}
% ===========================================================================

La versione scritta a mano processa 8~chiavi per iterazione. Due load \texttt{ymm} caricano 4+4~chiavi uint64. Ogni chiave viene XOR-folded shiftando a destra di 32~bit e facendo XOR con il valore originale, lasciando il risultato folded nei 32~bit bassi di ogni lane a 64~bit. Gli 8~valori folded vengono poi impacchettati in un singolo registro \texttt{ymm}: \texttt{permutevar8x32} con indici \texttt{[0,2,4,6,1,3,5,7]} estrae i valori low-32 da ciascun registro, e \texttt{permute2x128} con selettore \texttt{0x20} combina i 128~bit inferiori di entrambi in un unico registro di 8$\times$uint32. Un singolo \texttt{vpmulld} esegue tutte le 8~moltiplicazioni Fibonacci, un right-shift estrae i partition ID, e uno store a 256~bit scrive i risultati. Le restanti $N \bmod 8$ chiavi sono gestite da una coda scalare.

Il binario AVX2 verifica la correttezza rispetto a un riferimento scalare interno (compilato con \texttt{-fno-tree-vectorize} tramite pragma) confrontando i checksum FNV-1a prima del benchmarking; per $N\!\leq\!32$ stampa anche il confronto elemento per elemento.

% ===========================================================================
\section{Analisi delle prestazioni}
\label{sec:analysis}
% ===========================================================================

Tutti i benchmark sono stati eseguiti su node09 (AMD EPYC 7301, Zen~1, DDR4-2666) con allocazione SLURM esclusiva (\texttt{--partition=gpu-excl}), 11~ripetizioni, mediana riportata.

\subsection{Risultati}

La Tabella~\ref{tab:nsweep} riporta throughput e speedup su cinque dimensioni dell'input ($P\!=\!256$).

\begin{table}[h]
\centering
\resizebox{\columnwidth}{!}{%
\small
\begin{tabular}{@{}rrccrcc@{}}
\toprule
& \textbf{Baseline} & \multicolumn{2}{c}{\textbf{Autovec}} & \multicolumn{2}{c}{\textbf{AVX2}} \\
\cmidrule(lr){2-2}\cmidrule(lr){3-4}\cmidrule(lr){5-6}
$N$ & Mkeys/s & Mkeys/s & $S$ & Mkeys/s & $S$ \\
\midrule
$10^6$   & 929  & 1377 & 1,48$\times$ & 1276 & 1,37$\times$ \\
$10^7$   & 915  & 1274 & 1,39$\times$ & 1172 & 1,29$\times$ \\
$5{\times}10^7$  & 908  & 1314 & 1,45$\times$ & 1168 & 1,29$\times$ \\
$10^8$   & 918  & 1310 & 1,43$\times$ & 1200 & 1,31$\times$ \\
$2{\times}10^8$  & 908  & 1322 & 1,46$\times$ & 1152 & 1,30$\times$ \\
\bottomrule
\end{tabular}%
}
\caption{Throughput e speedup CPU rispetto alla baseline con $P\!=\!256$.}
\label{tab:nsweep}
\end{table}

Con $N\!=\!10^8$ il kernel CUDA da solo raggiunge 67\,320~Mkeys/s (808~GB/s, 81\% del picco HBM2 dell'A30), ma end-to-end inclusi i trasferimenti PCIe solo 1\,031~Mkeys/s ($1{,}12\times$); vedere Sezione~\ref{sec:cuda}.

\subsection{Perché lo speedup è limitato}

Ogni chiave richiede 12~byte di traffico di memoria (8\,B lettura + 4\,B scrittura) e solo $\sim$4 operazioni intere. L'intensità operazionale è $I = 4/12 \approx 0{,}33$ ops/byte, collocando il kernel nella regione bandwidth-bound del modello roofline. La banda single-core su node09 è stata misurata a $\sim$16,4~GB/s; la versione autovec raggiunge 15,7~GB/s, circa il 96\% di questo tetto (Figura~\ref{fig:roofline}). SIMD accelera la parte di calcolo, ma i trasferimenti di memoria --- che dominano --- restano invariati; lo speedup di 1,43$\times$ deriva da un utilizzo più efficiente della banda, non da un calcolo più veloce.

\begin{figure}[h]
\centering
\includegraphics[width=\columnwidth]{../results/plots/14_roofline.png}
\caption{Modello roofline. Tutte le varianti si trovano sulla rampa di banda ($I\!=\!0{,}33$), ben al di sotto dei tetti di calcolo.}
\label{fig:roofline}
\end{figure}

\subsection{Autovec vs.\ intrinsics AVX2}

In modo forse controintuitivo, la versione vettorizzata dal compilatore supera gli intrinsics scritti a mano (15,7 vs 14,4~GB/s). Poiché il kernel è limitato dalla memoria, le differenze marginali nello scheduling delle istruzioni, loop unrolling e allocazione dei registri che GCC applica durante l'auto-vettorizzazione si traducono in un utilizzo della banda leggermente migliore. Il codice con intrinsics impone una sequenza di istruzioni rigida che il compilatore non può facilmente riordinare.

\subsection{Sweep dei parametri}

\textbf{N-sweep} (Figura~\ref{fig:speedup}): lo speedup è stabile su tutte le dimensioni dell'input, da $10^6$ a $2{\times}10^8$ chiavi. Questo conferma che il kernel è bandwidth-bound indipendentemente da $N$.

\textbf{P-sweep} (Figura~\ref{fig:sweeps}a): throughput e speedup sono indipendenti da $P$ (tutte le curve piatte su $P \in [2, 1024]$), confermando che cambiare $P$ modifica solo la costante di shift --- traffico di memoria e calcolo per chiave restano invariati.

\textbf{Key-space sweep} (Figura~\ref{fig:sweeps}b): il throughput è invariante rispetto alla dimensione del key-space (tasso di duplicati), come atteso per una hash branchless a costo costante per chiave.

\begin{figure}[h]
\centering
\includegraphics[width=\columnwidth]{../results/plots/18_psweep_keyspace_combined.png}
\caption{(a)~Throughput vs.\ $P$ con $N\!=\!10^8$: piatto su tutti i valori di partizioni. (b)~Throughput vs.\ dimensione del key-space con $N\!=\!10^8$, $P\!=\!256$: invariante al tasso di duplicati.}
\label{fig:sweeps}
\end{figure}

\begin{figure}[h]
\centering
\includegraphics[width=\columnwidth]{../results/plots/03_speedup_vs_N.png}
\caption{Speedup vs.\ dimensione dell'input $N$ ($P\!=\!256$). Sia auto-vectorized che AVX2 intrinsics mantengono uno speedup stabile su tutti gli $N$, coerente con un kernel bandwidth-bound.}
\label{fig:speedup}
\end{figure}

\subsection{Stabilità delle misure}

Tutti i 55 punti di misura hanno coefficiente di variazione inferiore al 5\%. Il binario autovec è il più stabile (CoV medio 0,4\%), mentre gli intrinsics AVX2 mostrano varianza leggermente maggiore (CoV medio 3,5\%), probabilmente perché lo scheduling di GCC produce pattern di accesso alla memoria più deterministici della sequenza scritta a mano.

% ===========================================================================
\section{CUDA}
\label{sec:cuda}
% ===========================================================================

Il kernel GPU usa la stessa hash, un thread per chiave, 256~thread/blocco. La correttezza è verificata tramite confronto di checksum FNV-1a rispetto al riferimento CPU; per $N\!\leq\!32$ viene stampato anche il confronto elemento per elemento.

\begin{table}[h]
\centering\small
\begin{tabular}{@{}lrrr@{}}
\toprule
\textbf{Fase} & \textbf{Tempo (ms)} & \textbf{BW (GB/s)} & \textbf{\%} \\
\midrule
Trasferimento H$\to$D & 63,2 & 12,7 & 65\% \\
Kernel                &  1,5 & 808  &  2\% \\
Trasferimento D$\to$H & 32,3 & 12,4 & 33\% \\
\midrule
\textbf{End-to-end} & \textbf{97,0} & --- & 100\% \\
\bottomrule
\end{tabular}
\caption{Tempi CUDA con $N\!=\!10^8$, $P\!=\!256$.}
\label{tab:cuda}
\end{table}

Il kernel da solo raggiunge 808~GB/s (81\% del picco HBM2 dell'A30, pari a 993~GB/s). I trasferimenti PCIe pesano per il 98\% del tempo end-to-end, quindi scaricare questo kernel isolato sulla GPU non dà vantaggi rispetto alla CPU. La GPU converrebbe solo con dati già presenti sul device o con stadi successivi della pipeline anch'essi su GPU.

% ===========================================================================
\section{Conclusioni}
% ===========================================================================

Questo kernel è un classico workload memory-bound. Con $I\!=\!0{,}33$ ops/byte, le prestazioni sono dettate dalla banda DRAM indipendentemente dalla strategia SIMD. La versione autovec raggiunge il 96\% del tetto di banda single-core, lasciando praticamente nessun margine di miglioramento su singolo thread. La scelta di una hash a 32~bit (\texttt{vpmulld} nativo) rispetto a una a 64~bit è stata fondamentale: con una hash a 64~bit, AVX2 era in realtà \emph{più lento} dello scalare a causa dell'overhead della decomposizione \texttt{vpmuludq} in un regime limitato dalla banda.

\begin{thebibliography}{2}
\bibitem{knuth1998}
D.\,E.~Knuth, \textit{The Art of Computer Programming, Volume~3: Sorting and Searching}, 2nd ed. Addison-Wesley, 1998, \S\,6.4, pp.\,513--558.

\bibitem{ferragina2023}
P.~Ferragina, \textit{Pearls of Algorithm Engineering}. Cambridge University Press, 2023, Cap.\,8: The Dictionary Problem, pp.\,96--127.
\end{thebibliography}

\end{document}
